% !TeX program = xelatex
% Evidence-aligned manuscript draft. DCVR-v2 is a proposal, not an evaluated method.
\documentclass[journal]{IEEEtran}
\usepackage{fontspec}
\setmainfont{TeX Gyre Termes}
\usepackage{xeCJK}
\setCJKmainfont[BoldFont=SimHei]{SimSun}
\setCJKsansfont{Microsoft YaHei}
\setCJKmonofont{SimSun}
\usepackage{cite}
\usepackage{amsmath,amssymb,amsfonts}
\usepackage{graphicx,booktabs,array,multirow}
\usepackage{balance}
\usepackage{xcolor}
\usepackage[hidelinks]{hyperref}
\renewcommand{\abstractname}{摘\quad 要}
\renewcommand{\figurename}{图}
\renewcommand{\tablename}{表}
\renewcommand{\IEEEkeywordsname}{关键词}
\newcommand{\R}{\mathrm{R}}
\newcommand{\T}{\mathrm{T}}
\newcommand{\softmax}{\operatorname{softmax}}
\newcommand{\sigmoid}{\operatorname{sigmoid}}
\newcommand{\pending}{\textit{待验证}}
\begin{document}

\title{面向退化场景的 RGB-T 跟踪前景条件协同 Value 路由：设计与证据边界}
\author{\textit{作者与单位待补充}}
\maketitle

\begin{abstract}
RGB-T 跟踪中的跨模态互补受场景退化、局部遮挡和背景干扰共同影响。模板多头自注意力（MSHA）能够在模板表示间传递信息，但注意力权重本身并不表示来源 token 对当前目标定位的效用。本文据此研究一个更窄的问题：在不改变原生查询—键匹配、定位头和跟踪状态更新的条件下，能否根据目标前景证据，分别控制跨模态消息的方向预算和来源模板 token 的相对贡献。我们提出前景条件协同 Value 路由（DCVR-v2）的可检验设计：从模板—搜索匹配分布提取软匹配强度、熵、局部竞争差和跨模态一致性；利用首帧框生成的软前景权重聚合模态状态；再将方向级消息预算与前景内 token 重分配分开计算，并只作用于跨模态 Value 路径。现有实验资产只验证了较早的 DCVR-v1 和 Early-conditioned MSHA gate，不能作为 DCVR-v2 的性能证据。DCVR-v1 的 LasHeR full-244 记录为 AUC/PR/NPR 55.96/69.72/66.04，但训练未获得前景 mask、测试获得了前景 mask；其与受保护 M1 的结果来源和软件环境也不匹配，故不能据此作严格因果比较。固定状态和受控退化诊断同样未显示一致的定位改善。本文因此将 DCVR-v2 定位为待验证方法，并给出能够闭合“输入证据—消息重分配—定位读出—闭环轨迹”的最小实验协议，不宣称已取得性能提升。
\end{abstract}

\begin{IEEEkeywords}
RGB-T 单目标跟踪，跨模态融合，前景条件路由，Value 消息，退化鲁棒性。
\end{IEEEkeywords}

\section{引言}
单目标跟踪需要根据首帧目标框，在后续视频中持续估计目标位置。可见光（RGB）提供颜色与纹理信息，但会受低照度、曝光变化和运动模糊影响；热红外（TIR）较少依赖可见光照明，却可能受到热背景、低温差和边缘模糊影响。LasHeR 等基准推动了 RGB-T 跟踪方法与评价协议的发展\cite{li2019benchmark,li2022lasher}。双模态输入提供互补线索，但互补性不意味着每一帧、每个位置和每个模板 token 都具有相同可靠性。

模板—搜索交互和光照引导融合已被用于 RGB-T 跟踪\cite{hui2023tbsi,li2025tbsi}。然而，注意力分数描述的是特征相似关系，不等于“该来源 token 的 Value 是否有助于当前目标定位”。当来源模板包含背景或退化线索时，跨模态消息可能传播无关内容。另一方面，仅用全局模态权重难以表达模板内部的目标与背景差异；仅按搜索匹配尖锐度加权，也可能把稳定的背景峰误作目标证据。

本文的研究问题是：\emph{如何以首帧目标区域为条件，对模板 MSHA 的跨模态 Value 消息进行方向级与 token 级的协同分配，同时保持原有注意力匹配和跟踪读出路径？}我们将方向预算与 token 分配拆成两项：方向预算回答“从哪一模态向哪一模态传多少跨模态消息”，token 分配回答“该来源模态内部哪些前景模板 token 相对更值得传递”。将两者分开也避免了一个归一化陷阱：如果全局方向偏置加到 token 分数后再按前景均值归一化，该常数偏置会被完全消去。

该设计来自一条尚未闭合的证据链。受保护 M1 的固定参考为 LasHeR full-244 AUC/PR/NPR 56.21/70.21/66.20。M1 的失配分析提示部分严重失败始于目标几何仍位于搜索裁剪内，之后反馈漂移并导致目标出界；这说明跟踪错误具有闭环性质，但没有将首因归结为 MSHA 或融合层。Early-conditioned MSHA gate 的 full-244 点估计相对 M1 小幅上升，但其历史配对分析后来因逐序列预测来源不一致而失效。已实现的 DCVR-v1 也完成 full-244，但训练与测试前景 mask 不一致，且现存 M1 预测目录的重算值与受保护日志不符。因此这些结果能够界定风险和证据缺口，不能验证 DCVR-v2。

本文贡献限定为三点：
\begin{itemize}
  \item 将跨模态 Value 路由拆解为方向预算与前景条件 token 分配，给出可直接实现和证伪的 DCVR-v2 数学定义。
  \item 根据已有固定状态、受控退化、代理训练和 full-244 资产，区分已观察事实、不可归因结果与尚未验证的机制，不把模块活动性等同于跟踪效用。
  \item 给出按相同初始化、相同数据合同和独立输出目录执行的实验闭环，覆盖输入合同、消息变化、CENTER 读出、闭环指标、效率及失败归因。
\end{itemize}
本文不声称 DCVR-v2 已实现、已超过 M1，或已证明可以缓解特定退化类型。

\section{相关工作与问题定位}
\subsection{RGB-T 跟踪与模板交互}
RGB-T 基准和跟踪器研究了双模态特征提取、融合和在线定位\cite{li2019benchmark,li2022lasher,tang2025mvrgbt}。TBSI 类方法通过模板桥接跨模态信息，并结合搜索区域融合与模板策略\cite{hui2023tbsi,li2025tbsi}。本文沿用其模板—搜索跟踪框架，只考察模板 MSHA 中跨模态 Value 的来源分配；不改变在线模板更新，不替换 CENTER 定位头，也不把固定状态诊断当作闭环跟踪结果。

\subsection{模态质量与动态路由}
质量感知融合、模板不确定性和动态路由方法从不同角度处理模态差异\cite{ding2025qstnet,ding2025tumnet,ding2026routing}。本文的差异不在于笼统提出“动态权重”，而在于将可观测的模板前景区域作为条件，并把方向预算与来源 token 分配显式分解。该分解仍属于注意力 Value 路径上的条件重加权；在完成系统检索和直接比较前，本文不主张其为新的注意力原理，也不声称首次使用目标条件融合。

\subsection{研究假设及可证伪边界}
DCVR-v2 的工作假设是：在部分目标仍处于搜索裁剪内的失配状态中，不同模态的前景模板 token 对跨模态消息具有不同效用；目标条件的相对分配可能比整路缩放更适合表达这种差异。若正确首帧前景 mask 不优于等面积移位 mask，若预算与 token 权重改变却不改变 CENTER 排名或解码框，或若 matched 训练比较不能改善闭环结果，则该假设不获支持。即使局部消息和读出改善，也不能自动推出全量数据或多 seed 提升。

\section{方法}
\subsection{跟踪器与符号}
令 $m\in\{\R,\T\}$ 表示模态，$T_m\in\mathbb{R}^{M\times d}$ 和 $S_m\in\mathbb{R}^{N\times d}$ 分别为同模态模板与搜索 token。首帧标注框映射为模板软前景权重 $w_i^m\in[0,1]$。原生 MSHA 使用查询—键分数生成注意力，并在模板 token 间聚合 Value；DCVR-v2 保留该匹配分数，只调节跨模态来源 Value 的权重。对接收模态 $n$，其来自模态 $m$ 的跨模态消息记为 $M_{n\leftarrow m}$，原生同模态消息保持不变。

\subsection{模板条件的匹配分布证据}
对来源模板 token $i$，先计算它与同模态搜索 token $j$ 的归一化相似度：
\begin{equation}
 s_{ij}^{m}=\frac{\langle \operatorname{norm}(W_T^mT_{m,i}),
                   \operatorname{norm}(W_S^mS_{m,j})\rangle}{\tau_s},
 \qquad p_{ij}^{m}=\softmax_j(s_{ij}^{m}).
\end{equation}
从分布而非单一最大值提取四类证据：
\begin{align}
 p_i^m &= \sum_j p_{ij}^m s_{ij}^m, \\
 h_i^m &= -\frac{1}{\log N}\sum_j p_{ij}^m\log(p_{ij}^m+\epsilon), \\
 c_i^m &= \max_{j\in\mathcal{N}_i}s_{ij}^m-\max_{j\notin\mathcal{N}_i}s_{ij}^m, \\
 a_i^m &= \operatorname{cos}(T_{m,i},T_{\bar m,\pi(i)}).
\end{align}
其中 $j_i^*=\arg\max_j s_{ij}^m$，$\mathcal{N}_i$ 是以 $j_i^*$ 为中心、尺度由首帧目标框确定的搜索邻域，$\pi(i)$ 是两模态模板 token 的空间对应索引。$p_i^m$ 表示匹配强度，$h_i^m$ 表示分布不确定性，$c_i^m$ 比较局部峰与非邻域峰，$a_i^m$ 表示对应模板 token 的跨模态一致性。它们只是路由输入，不是校准概率，也不天然等于目标定位效用。令 $e_i^m=[p_i^m,h_i^m,c_i^m,a_i^m]$。

\subsection{前景条件模态状态}
利用首帧框得到的软权重计算模态级状态：
\begin{equation}
 q_m=\frac{\sum_i w_i^m\,\phi(e_i^m)}
           {\sum_i w_i^m+\epsilon},\qquad
 \bar a_{\R\T}=\frac{\sum_i w_i^\R w_{\pi(i)}^\T a_i^\R}
 {\sum_i w_i^\R w_{\pi(i)}^\T+\epsilon}.
\end{equation}
其中 $\phi$ 是共享的小型映射。背景模板 token 不直接支配模态级状态；软边界可降低目标框与 patch 网格不对齐造成的硬切分误差。前景 mask 的生成与使用必须在训练和推理中一致。缺失 mask 时不得静默使用全一权重；应在数据入口显式构造，或将该样本标为协议失败。

\subsection{方向预算与 token 级分配}
对每个有向跨模态通路，方向网络读取两模态状态、跨模态一致性及状态差：
\begin{equation}
 d_{m\rightarrow n}=f_{\mathrm{dir}}([q_m,q_n,\bar a_{mn},q_m-q_n]).
\end{equation}
对两个相反的跨模态方向进行归一化，得到有界方向预算：
\begin{equation}
 \begin{bmatrix}b_{\R\rightarrow\T}\\b_{\T\rightarrow\R}\end{bmatrix}
 =2\,\softmax\!\left(\begin{bmatrix}d_{\R\rightarrow\T}\\
 d_{\T\rightarrow\R}\end{bmatrix}\right),\qquad
 b_{\R\rightarrow\T}+b_{\T\rightarrow\R}=2.
\end{equation}
因而一个方向可相对增强时，另一个方向的预算相应减少；两路平均预算为 1。该预算只作用于跨模态支路，不缩放同模态 Value。

来源 token 网络逐 token 读取本模态与接收模态证据：
\begin{equation}
 z_i^{m\rightarrow n}=f_{\mathrm{tok}}([e_i^m,e_{\pi(i)}^n,w_i^m,a_i^m]).
\end{equation}
随后在前景区域内中心化并归一化：
\begin{align}
 \tilde z_i^{m\rightarrow n}&=z_i^{m\rightarrow n} -
 \frac{\sum_k w_k^m z_k^{m\rightarrow n}}{\sum_k w_k^m+\epsilon},\\
 g_i^{m\rightarrow n}&=\exp\{\beta\tanh(\tilde z_i^{m\rightarrow n})\},\\
 \hat g_i^{m\rightarrow n}&=\frac{g_i^{m\rightarrow n}}
 {\frac{\sum_k w_k^m g_k^{m\rightarrow n}}{\sum_k w_k^m}}.
\end{align}
因此前景区域内 $\hat g$ 的加权均值为 1，$\beta$ 控制 token 间重分配幅度。方向预算没有被放入这一步的均值归一化，避免全局方向项被代数抵消。

\subsection{跨模态 Value 更新与恒等初始化}
原生 MSHA 对拼接后的 RGB/TIR 模板 token 做联合注意力。将其输出按 query 模态和 key 来源模态分解，记 $A_{n\leftarrow m}$ 为接收模态 $n$ 的 query 到来源模态 $m$ 模板 token 的注意力子矩阵，$V_m$ 为对应 Value。DCVR-v2 只重加权跨模态项，同模态项保持原样：
\begin{equation}
 M_{n\leftarrow m}=A_{n\leftarrow m}
 \left[(b_{m\rightarrow n}\hat g_i^{m\rightarrow n})\odot V_{m,i}\right].
\end{equation}
所有同模态消息、Q/K 分数、原有残差和 CENTER 头保持不变。方向网络与 token 网络最后一层零初始化，使初始 $d=0$、$z=0$，从而 $b=1$、$\hat g=1$，候选输出与 parent 精确一致；零初始化投影仍可获得梯度。训练后由原始跟踪损失学习。零初始化、参数保存加载、梯度路径和数值稳定性是实现合同检查，不是性能证据。

\subsection{优化目标与限制}
初版不引入人工退化标签或额外质量监督，沿用基线的分类、框回归和质量损失：
\begin{equation}
 \mathcal{L}=\mathcal{L}_{\mathrm{base}}.
\end{equation}
因此路由是否与定位效用一致必须由独立诊断检验，不能从训练 loss 下降推断。该设计无法在两种模态都缺失目标信息时创造线索，也不处理重检测、长期记忆和目标出 crop 后恢复。实现还需记录路由极值、逐 token 权重、梯度、显存与延迟，防止均值统计掩盖少数 token 的异常放大。

\section{实验协议与证据闭环}
\subsection{锚点、数据和指标}
受保护 parent 为 M1：\texttt{TBSI\_E\_MHSA\_FULL\_4090D\_seed42}，配置与 checkpoint 及其固定报告按实验资产数据库登记。其 LasHeR full-244 参考 AUC/PR/NPR/OP50/OP75 为 56.21/70.21/66.20/68.33/47.13，环境为 1$\times$ RTX 4090D、Python 3.8.20、PyTorch 1.9.0+cu111、CUDA 11.1、seed 42。该值是受保护记录，不用当前被改写的 M1 预测目录重新生成配对差值。

新候选必须有独立配置、输出目录、checkpoint 和资产 ID。训练/测试数据范围、manifest、预处理、更新策略、随机种子、评测脚本、权重和逐序列预测 hash 均须保存。主指标为 LasHeR AUC、PR、NPR、OP50、OP75；另报告逐序列配对差异、置信区间、胜负数、推理吞吐、参数量和延迟。full-244 单 seed 仅支持该次运行的描述性结论，不等于多 seed 稳健性。

\subsection{分阶段验证}
完整闭环包含下列预注册阶段，前一阶段失败时停止升级，不将其结果包装成完整验证：
\begin{enumerate}
  \item \textbf{实现合同：} disabled parity、mask 训练/推理一致、shape、零初始化、保存加载、有限值、单 batch 梯度和冻结参数状态检查。
  \item \textbf{固定状态机制：} 在不与信号拟合重叠的 train-side 序列上，逐帧记录前景/背景 token 权重、两方向消息预算、Value 改变量、CENTER 加窗前后排名、偏移/尺度解码和同 crop IoU。比较正确 mask、等面积移位 mask、均匀 token 权重及无方向路由。
  \item \textbf{受控退化：} 对 RGB 低照度/模糊、TIR 对比度/噪声及局部遮挡进行成对干预；检验退化来源方向是否受抑、另一方向是否得到相对预算，以及干净条件是否受损。合成干预只支持所测变换范围内的机制判断。
  \item \textbf{匹配代理训练：} 从相同 SOT 初始 state 克隆公共参数与新参数初值，使用相同 manifest、seed、优化器、步数和 worker 合同，只改变 DCVR-v2 开关。与继续训练 control 成对比较。
  \item \textbf{闭环 full-244：} 仅当合同、固定状态和匹配代理满足预注册门槛后执行。全量测试固定权重、评测口径和 tracker 参数；保留逐序列结果、错误片段和恢复轨迹。
  \item \textbf{消融及效率：} 逐项移除前景条件、分布熵/竞争差、一致性、方向预算或 token 分配；另设等范数缩放控制。报告额外参数、显存、延迟和 FPS，不把含诊断写盘的速度当成纯推理速度。
\end{enumerate}

\section{已有实验结果与解释边界}
\subsection{固定参考和先导候选}
表~\ref{tab:existing} 汇总目前可核对的数字。除受保护 M1 固定参考外，候选数值来自各自登记运行。训练起点、比较控制和结果来源并不完全相同，表中不计算统一的因果差值。DCVR-v2 尚无实验行。
\begin{table*}[t]
\centering
\caption{已有记录的描述性结果。不同训练合同或结果来源间不作严格配对推断；DCVR-v2 尚未实现和评测。}
\label{tab:existing}
\small
\begin{tabular}{lccccccp{6.1cm}}
\toprule
运行 & AUC & PR & NPR & OP50 & OP75 & FPS & 证据状态与限制\\
\midrule
M1 受保护固定参考 & 56.21 & 70.21 & 66.20 & 68.33 & 47.13 & NA & full-244 固定记录；1$\times$ RTX 4090D，PyTorch 1.9.0+cu111。\\
Early-conditioned MSHA gate & 56.56 & 70.58 & 66.76 & 68.69 & 47.64 & 53.14 & full-244 单 seed 点估计；历史配对分析 provenance invalidated，联合训练且无 matched gate-off full15 control。\\
DCVR-v1 full15 & 55.96 & 69.72 & 66.04 & 68.06 & 47.09 & 9.73 & full-244 已完成；训练前景 mask fallback 为 100\%，测试为 0\%；与 M1 的目录复算及软件版本不一致。FPS 含诊断写盘且测试线程条件变化。\\
DCVR-v2 & -- & -- & -- & -- & -- & -- & 仅为本文定义的后续候选；没有训练、推理或闭环结果。\\
\bottomrule
\end{tabular}
\end{table*}

Early gate 的已登记 full-244 点估计均略高于 M1，但逐序列来源审计发现 parent 与 candidate 各有 7 条预测不匹配，当前 M1 目录重算 AUC 56.766，不符合受保护 M1 日志的 56.21。另一个 fixed-state A/B 中 Hann top-1 变化为 0/28，未观察到中心排名改善。故允许的结论是“曾记录方向性正点估计，但性能增益和机制未确立”。

DCVR-v1 的固定状态测试观察到路由非恒等，但 IoU 变化为 $-0.00062$，CENTER 排名仅 7 个事件中 1 个改变。受控退化 smoke 中，双模态 blur、contrast、lowlight 条件的 IoU 差分别为 $-0.00079$、$+0.00068$、$+0.00169$，clean 条件为 $+0.00554$；方向依赖条件，样本规模仅 8 个事件。quick3 候选相对控制的 AO 均值差为 $+0.04726$，主要由单一序列贡献，不能作为稳定增益证据。full15 DCVR-v1 汇总值可复算，但训练未获前景 mask、测试获得前景 mask；当前 M1 预测目录复算与 M1 保护日志不符，且 PyTorch/CUDA 版本不同。因此既不能把点估计下降单独归因于算子，也不能用现有配对 bootstrap 修复这些合同差异。

\subsection{允许的结论与未决解释}
现有数据支持如下有限判断：
\begin{itemize}
  \item \textbf{已观察：} Early gate 与 DCVR-v1 能改变消息路径；DCVR-v1 full-244 已完成，汇总指标低于受保护 M1 的登记点估计。
  \item \textbf{暂不支持：} Early gate 或 DCVR-v1 已稳定改善 M1；当前路由信号已准确估计目标定位效用；受控退化结果能推广到真实退化分布。
  \item \textbf{尚未区分：} DCVR-v1 结果受前景 mask 训练/测试不一致、路由信号与目标效用错位、跨软件结果来源差异及单 seed 影响的相对程度。
  \item \textbf{待检验：} DCVR-v2 的分布证据、前景条件、方向预算和 token 分配是否能依次改变目标相关消息、中心读出及闭环轨迹。
\end{itemize}

\subsection{消融与报告矩阵}
表~\ref{tab:closure} 是需要完成的因果闭环，而非已有结果。每项必须有独立资产登记；不满足比较合同的运行只列为诊断或探索性结果。
\begin{table}[t]
\centering
\caption{DCVR-v2 的最小证据闭环状态。}
\label{tab:closure}
\small
\begin{tabular}{p{2.1cm}p{4.7cm}p{1.1cm}}
\toprule
证据层 & 必需比较与读数 & 状态\\
\midrule
输入合同 & 正确 mask 与训练/测试一致性；hash、预处理、checkpoint & 待完成\\
消息机制 & 方向预算、前景/背景分配、Value 改变量与稳定性 & 待完成\\
局部定位 & CENTER rank、margin、offset/size、同 crop IoU & 待完成\\
受控退化 & 模态单独及联合退化、clean control & 待完成\\
训练贡献 & 同初值、同数据和优化预算的 on/off 配对训练 & 待完成\\
闭环结果 & LasHeR full-244、逐序列 CI、失败轨迹和恢复统计 & 待完成\\
组件归因 & 前景、分布证据、方向预算、token 权重、等范数消融 & 待完成\\
成本与复现 & FPS、延迟、显存、参数、环境、完整命令与产物 hash & 待完成\\
\bottomrule
\end{tabular}
\end{table}

\section{讨论}
DCVR-v2 的可检验价值在于把“消息来源的相对可靠性”拆成方向与 token 两个尺度，而不是给 gate 堆叠更多输入。熵低可能意味着目标匹配集中，也可能意味着背景纹理造成稳定错误峰；跨模态一致性高可能代表互补，也可能代表共同背景。因此任何路由可视化都必须与 CENTER 响应和最终解码框一起检查。若消息权重变化没有传递到定位读出，应停止用机制图宣称有效。

当前 full-244 证据还有比较来源问题。M1 的受保护日志是唯一固定参照；若需要新候选的序列级因果比较，必须在独立输出目录中用受保护 checkpoint 和候选 checkpoint、同一数据 manifest、同一评测器及同一软件环境重新生成结果，并验证每条序列预测 hash。不得覆盖 M1 记录或复用其输出目录。单 seed 结果只能报告为单次观察；如果资源预算不足以完成多 seed，应相应收窄结论。

此外，现有机制诊断指出的系统级风险链为“crop 内定位失配—反馈漂移—目标出界—恢复不足”。DCVR-v2 直接作用于模板 Value 分配，不直接修改选峰、框解码、状态提交或重检测。因此即使前两层证据成立，仍需展示其确实减少闭环失配或缩短失败持续时间；否则应将结论限定为局部表示机制，不声称修复完整失败链。

\section{结论}
本文将 RGB-T 跟踪的跨模态消息分配表述为一个前景条件的方向预算与 token 级重分配问题，并给出 DCVR-v2 的可证伪设计。现有资产没有验证该版本。已完成的 Early gate 与 DCVR-v1 实验提供了先导信号，也暴露出定位读出未改善、训练/测试 mask 不一致和结果来源不可严格配对等限制。当前最强结论是：跨模态 Value 可以被调制，但其目标效用和闭环收益尚未建立。只有通过合同一致的机制诊断、matched 训练对照和独立 full-244 闭环评测后，才能判断 DCVR-v2 是否构成有效跟踪改进。

\balance
\bibliographystyle{IEEEtran}
\bibliography{references}
\end{document}
